\section*{Overview}

A trie stores strings by their prefixes. Instead of comparing one full string against many others, I walk character by
character through shared nodes. That makes tries a natural fit for dictionary problems, prefix queries, and some
automaton constructions.

\section*{When to Use It}

Use a trie when:

\begin{itemize}
  \item the main queries are by prefix,
  \item the alphabet is small enough to store transitions directly,
  \item many strings share prefixes and you want to reuse that structure.
\end{itemize}

\section*{Core Idea}

Each edge corresponds to one character, and each root-to-node path represents a prefix. Full words are marked with an
end flag or a terminal counter.

Because the path length equals the string length, operations depend on the string size rather than the number of stored
strings.

\section*{Key Insight}

The trie is valuable when the query is ``what happens after this prefix?'' If the problem is only about exact whole-word
lookup, a hash table is usually simpler. The structure pays for itself when prefix sharing matters.

\section*{Operations / Main Technique}

\begin{itemize}
  \item insert a word,
  \item test whether a word exists,
  \item test whether a prefix exists,
  \item count how many words pass through a prefix node.
\end{itemize}

\paragraph{Worked example.}
If the set contains \(\texttt{car}\), \(\texttt{card}\), and \(\texttt{care}\), then the path for \(\texttt{car}\) is
shared and branches only afterward.

\section*{Correctness Intuition}

Each inserted character follows or creates the unique outgoing edge for that symbol, so the path stored in the trie is
exactly the word itself. Shared prefixes meet at shared nodes, and different next characters diverge into different
children.

\section*{Complexity Analysis}

For word length \(L\):
\begin{itemize}
  \item insert: \(O(L)\),
  \item contains: \(O(L)\),
  \item prefix query: \(O(L)\).
\end{itemize}

\section*{Implementation}

The sample code uses a lowercase English alphabet trie with:

\begin{itemize}
  \item a terminal counter,
  \item a pass-through counter,
  \item insert, exact lookup, and prefix counting.
\end{itemize}

\section*{Common Pitfalls}

\begin{itemize}
  \item Using a trie when the alphabet is huge and sparse.
  \item Forgetting whether repeated insertions should be allowed or counted.
  \item Confusing ``prefix exists'' with ``full word exists''.
  \item Spending too much memory on a static array when a compressed representation would be better.
\end{itemize}

\section*{Variants / Extensions}

\begin{itemize}
  \item Binary trie for XOR problems.
  \item Aho-Corasick, which augments the trie with failure links.
  \item Compressed trie or radix tree for memory-sensitive settings.
\end{itemize}

\section*{Practice Problems}

\begin{itemize}
  \item Prefix dictionary queries.
  \item Count how many stored words share a given prefix.
  \item Maximum XOR using a binary trie.
\end{itemize}

\section*{References}

\begin{itemize}
  \item \href{https://cp-algorithms.com/string/aho_corasick.html}{cp-algorithms: Trie-based Aho-Corasick background}.
  \item \href{https://usaco.guide/gold/strings}{USACO Guide: String Techniques}.
  \item \href{https://codeforces.com/catalog}{Codeforces Catalog}.
\end{itemize}
